

\section{Basics}

\begin{frame}{Qu'est-ce que \mongodb/,?}

%\frametitle{Qu'est-ce que \mongodb/,?}

\mongodb/, un système "NoSQL", l'un des plus populaires

\begin{redbullet}
  \item s'appuie sur un modèle de données semi-structuré (encodage JSON)\,;
  \item pas de schéma (complète flexibilité)\,;
  \item un \alert{langage d'interrogation} original, limité (et spécifique)\,;
   \end{redbullet}
   \vfill
   
Construit dès l'origine comme un système \alert{scalable} et \alert{distribué}.

\begin{redbullet}
 \item distribution par partionnement (\textit{sharding})\,;
  \item tolérance aux pannes par réplication.
\end{redbullet}   


\vfill

Un aperçu du \alert{langage d'interrogation} de \mongodb/.


\end{frame}


\begin{frame}
\frametitle{Découvrons (avec le client \texttt{Compass})}

Récupérez Compass et connectez-vous à MongoDB Atlas
\vfill

Choisissez la base Patients et faite quelques essais

\vfill

\begin{redbullet}
  \item Insérez un document
  \item Lancez un \texttt{find} avec une requête vide
  \item Accédez au \texttt{shell} Mongo
\end{redbullet}

\vfill

\begin{noteblock}
Si vous voulez votre propre base vous pouvez récupérer les jeux
de données sur le site.
\end{noteblock}
\end{frame}




\begin{frame}[fragile]
\frametitle{Aperçu du langage de requêtes}

L'équivalent du \texttt{select-from-where}, \alert{sur une
seule collection} (pas de \alert{jointure}).

\vfill

Principe de \textit{pattern matching}\,: on apparie
les documents à la requête.


\begin{smallverbatim}
find ({"_id": "100"})
\end{smallverbatim}


\vfill

\begin{smallverbatim}
find ({"prenom": "aaron"})
\end{smallverbatim}

\vfill

Avec des \alert{expressions de chemin} (simples!)

\vfill


\begin{smallverbatim}
find ({"patholohies.intitule": "Embolie pulmonaire"} )
\end{smallverbatim}

\end{frame}


\section{Manipulation de données avec \mongodb/}


\begin{frame}[fragile]
\frametitle{Les requêtes}

Recherche par clé (très rapide). 

\begin{minted}[baselinestretch=.9,
               fontsize=\small,
               framesep=2mm]{javascript}
find({_id: "200"})
\end{minted}

\vfill 
\begin{minted}[baselinestretch=.9,
               fontsize=\small,
               framesep=2mm]{javascript}
find({"pathologies.dept": "57"})
\end{minted}

Avec pagination
\begin{minted}[baselinestretch=.9,
               fontsize=\small,
               framesep=2mm]{javascript}
db.patients.find({"sexe": "f"}).skip(0).limit(30)
\end{minted}

Avec une \alert{projection}

\begin{minted}[baselinestretch=.9,
               fontsize=\small,
               framesep=2mm]{javascript}
db.patients.find({"pathologies.dept": "57"}, {"prenom": 1})
\end{minted}

Avec une expression régulière (pas de guillemets!)\,:

\begin{minted}[baselinestretch=.9,
               fontsize=\small,
               framesep=2mm]{javascript}
db.patients.find({"prenom":  /aba/})
\end{minted}

Par intervalle:

\begin{minted}[baselinestretch=.9,
               fontsize=\small,
               framesep=2mm]{javascript}
db.patients.find({"pathologies.dept": {$gte: "50", $lte: "60"}})
\end{minted}

\end{frame}



\begin{frame}[fragile]
\frametitle{Les requêtes, suite}

Quelques clauses ensemblistes 

\begin{verbatim}
db.patients.find( { prenom: { \$in: ['aaron', 'adam', 'adeline' ] } } )
\end{verbatim}

Autres possibilités\,: \texttt{\$nin}, \texttt{\$all}, \texttt{\$exist}

\begin{verbatim}
db.patients.find( { prenom: { \$nin: ['aaron', 'adam', 'adeline' ] } } )
\end{verbatim}

\vfill

\begin{noteblock}
Les agrégations existent, mis c'est très compliqué
\end{noteblock}
\end{frame}




\begin{frame}[fragile]
\frametitle{Et les jointures}

Pas de jointure avec Mongo.
Il faut les effectuer du \alert{côté client}. 

\vfill

\figSlideWithSize{jointure-serveur-client}{7cm}

\vfill

En règle générale, il faut programme côté applicatif tout ce qu'on a perdu de SQL
\end{frame}

